% !TeX root = ../../Book4.tex
% 纲要标题：### 9.4 一次 Ext、Tor 与正合性判据
% 纲要标签：b4:sec:0804
% 纲要位置：计划纲要.md，第 3420 行。

\section{一次 Ext、Tor 与正合性判据}
\label{b4:sec:0804}

一次 Ext 或 Tor 的消失能检测投射性、内射性或平坦性，但测试对象必须遍历相应的全部模。整数模与多项式理想的例子可以把正合性的失败具体算出来。


% 纲要内容（仅作写作计划，尚非教材正文）：
%
% * Ext 检验投射性和内射性
% * Tor 检验平坦性
% * 回到第二卷反例解释失败机制
%

\subsection{Ext 与投射性、内射性}
\label{b4:subsec:0804:01}

\begin{theorem}\label{b4:thm:ext-projective-injective}
设 \(R\) 是含幺环，\(P,I\) 是幺性左 \(R\) 模。下列判据成立：
\begin{enumerate}
\item \(P\) 投射，当且仅当对所有幺性左 \(R\) 模 \(N\)，\(\operatorname{Ext}_R^1(P,N)=0\)；这也等价于对所有这样的 \(N\) 及所有 \(n>0\)，\(\operatorname{Ext}_R^n(P,N)=0\)。
\item \(I\) 内射，当且仅当对所有幺性左 \(R\) 模 \(M\)，\(\operatorname{Ext}_R^1(M,I)=0\)；这也等价于对所有这样的 \(M\) 及所有 \(n>0\)，\(\operatorname{Ext}_R^n(M,I)=0\)。
\end{enumerate}
\end{theorem}
\begin{proof}
投射或内射时用集中于零次的相应消解，正次 Ext 消失。反过来，设对所有左模 \(N\) 都有 \(\operatorname{Ext}_R^1(P,N)=0\)。任取满射 \(q:B\to C\)，令 \(A=\ker q\)，第二变量长正合列给出
\[
\operatorname{Hom}_R(P,B)\xrightarrow{q\circ-}\operatorname{Hom}_R(P,C)
\longrightarrow\operatorname{Ext}_R^1(P,A)=0.
\]
所以每个 \(P\to C\) 都是某个 \(P\to B\) 与 \(q\) 的复合，即 \(P\) 有投射提升性质。

同样，设对所有左模 \(M\) 都有 \(\operatorname{Ext}_R^1(M,I)=0\)。任取单射 \(j:A\to B\)，令 \(C=\operatorname{coker}j\)，第一变量长正合列给出
\[
\operatorname{Hom}_R(B,I)\xrightarrow{-\circ j}\operatorname{Hom}_R(A,I)
\longrightarrow\operatorname{Ext}_R^1(C,I)=0.
\]
所以每个 \(A\to I\) 都能延拓为 \(B\to I\)，即 \(I\) 内射。
\end{proof}

“对所有模”的量词有实质作用。例如 \(\operatorname{Ext}_{\mathbb Z}^1(\mathbb Z/2\mathbb Z,\mathbb Q)=0\)，并不能推出 \(\mathbb Z/2\mathbb Z\) 投射：这次消失来自第二变量 \(\mathbb Q\) 的内射性，而换成第二变量 \(\mathbb Z\) 时，Ext 就是 \(\mathbb Z/2\mathbb Z\)。又如 \(\operatorname{Ext}_{\mathbb Z}^1(\mathbb Z,\mathbb Z)=0\)，不能推出第二个 \(\mathbb Z\) 内射：这次消失来自第一变量 \(\mathbb Z\) 的投射性，换成第一变量 \(\mathbb Z/2\mathbb Z\) 就得到同一个非零 Ext。检验投射性时固定第一变量、遍历第二变量；检验内射性时则相反。

\subsection{Tor 与平坦性}
\label{b4:subsec:0804:02}

\begin{theorem}\label{b4:thm:tor-flat-criterion}
设 \(R\) 是含幺环，\(N\) 是幺性左 \(R\) 模。下列条件等价：\(N\) 平坦；对每个幺性右 \(R\) 模 \(M\)，\(\operatorname{Tor}_1^R(M,N)=0\)；对每个这样的 \(M\) 及每个 \(n>0\)，\(\operatorname{Tor}_n^R(M,N)=0\)。右模的判据以相反侧的幺性左模作测试。
\end{theorem}
\begin{proof}
若 \(N\) 平坦，\(-\otimes_RN\) 正合，故右模投射消解在张量后仍在正次数正合；于是所有正次 Tor 为零。全消失当然蕴含一次消失。若一次均消失，对任意右模单射 \(A\to B\)，令 \(C=B/A\)。长正合列的片段
\[
\operatorname{Tor}_1^R(C,N)\longrightarrow A\otimes_RN
\longrightarrow B\otimes_RN
\]
表明后一个箭头单射。对任意右模短正合列 \(0\to A\to B\to C\to0\)，张量积的右正合性已经保证 \(A\otimes_RN\to B\otimes_RN\to C\otimes_RN\to0\) 正合；缺少的只是最左边的单射性。现在这一步也成立，故 \(-\otimes_RN\) 正合，即 \(N\) 平坦。反环给出右模版本。
\end{proof}

不必为了判定平坦性计算所有高次 Tor；一次已经足够。相反，只对某一个模计算一次 Tor，通常只能发现反例，不能得出整体平坦性。

\subsection{非正合性的模论反例}
\label{b4:subsec:0804:03}

设 \(m>1\) 为整数。把 \(0\to\mathbb Z\xrightarrow{\times m}\mathbb Z\to\mathbb Z/m\mathbb Z\to0\) 与 \(\mathbb Z/m\mathbb Z\) 作张量，乘 \(m\) 的映射变为零。由于 \(\mathbb Z\) 投射，相应长正合列的片段是
\[
0\longrightarrow\operatorname{Tor}_1^{\mathbb Z}(\mathbb Z/m\mathbb Z,\mathbb Z/m\mathbb Z)
\longrightarrow\mathbb Z/m\mathbb Z\xrightarrow{0}\mathbb Z/m\mathbb Z.
\]
零映射的核是整个 \(\mathbb Z/m\mathbb Z\)，故 \(\operatorname{Tor}_1^{\mathbb Z}(\mathbb Z/m\mathbb Z,\mathbb Z/m\mathbb Z)\cong\mathbb Z/m\mathbb Z\)，它记录了张量后丢失的单射性。对同一列施加 \(\operatorname{Hom}_{\mathbb Z}(-,\mathbb Z)\)，得到
\[
\mathbb Z\xrightarrow{\times m}\mathbb Z
\longrightarrow\operatorname{Ext}_{\mathbb Z}^1(\mathbb Z/m\mathbb Z,\mathbb Z)
\longrightarrow0.
\]
乘 \(m\) 的映射不是满射，其余核就是 \(\operatorname{Ext}_{\mathbb Z}^1(\mathbb Z/m\mathbb Z,\mathbb Z)\cong\mathbb Z/m\mathbb Z\)。对同一个整数短正合列，张量积丢失的单射性由 Tor 的核项记录，反变 Hom 丢失的满射性由 Ext 的余核项记录。

\ProseParagraphBreak
再取域 \(k\) 上的多项式环 \(R=k[x,y]\) 及其理想 \(\mathfrak m=(x,y)\)，将 \(k\) 识别为 \(R/\mathfrak m\)。因为 \(R\) 交换，以下将这些模同时视为右 \(R\) 模。短正合列 \(0\to\mathfrak m\to R\to k\to0\) 在第一变量的 Tor 长正合列中给出
\[
\underbrace{\operatorname{Tor}_2^R(R,k)}_{=0}\longrightarrow
\operatorname{Tor}_2^R(k,k)\longrightarrow\operatorname{Tor}_1^R(\mathfrak m,k)
\longrightarrow\underbrace{\operatorname{Tor}_1^R(R,k)}_{=0}.
\]
两端因 \(R\) 投射而消失，再代入上一节算出的二次 Tor，得到
\[
\operatorname{Tor}_1^R(\mathfrak m,k)\cong
\operatorname{Tor}_2^R(k,k)\cong k.
\]
故 \(\mathfrak m\) 不平坦，尽管它是整环中的无挠子模。这具体说明主理想整环上的无挠模平坦这一结论不能直接推广到任意整环。上面的连接态射已经将二次 Tor 与核模的一次 Tor 对应；下一节说明何时能这样降低次数。
